\ifdefined\gkver\else\def\gkver{student}\fi
\documentclass[\gkver]{gaokaozhenti}
\begin{document}

\chapter{2015年安徽理综}
%% paperType: 理综物理部分

\begin{enumerate}
\item
%% number: 14
%% paperName: 2015年安徽理综第 14 题 6 分
%% typeId: 1
%% type: 单选题
%% chapter: 电场
%% point: 库仑定律
%% method: 无
%% score: 6
%% degree: 650
%% duplicateId: 0
%% body:
图示是$\alpha$粒子（氦原子核）被重金属原子核散射的运动轨迹，M、N、P、Q 是轨迹上的四点，在散射过程中可以认为重金属原子核静止不动。图中所标出的 $\alpha$ 粒子在各点处的加速度方向正确的是\xzanswer{C}

\onepicture[h=2.4cm]{figs/14.png}

\fourchoices[answer=C]
{M 点}
{N 点}
{P 点}
{Q 点}

%% answer: C
%% memo:
\memoanswer{
	由牛顿第二定律知加速度的方向与所受合力方向相同，$\alpha$ 粒子与重金属原子核的作用力为斥力，沿两连线的背离方向，由题图可得 P 点描述正确。选项 C 正确。
}

\item
%% number: 15
%% paperName: 2015年安徽理综第 15 题 6 分
%% typeId: 1
%% type: 单选题
%% chapter: 电场
%% point: 库仑定律
%% method: 无
%% score: 6
%% degree: 650
%% duplicateId: 0
%% body:
由库仑定律可知，真空中两个静止的点电荷，带电量分别 $q_{1}$ 和 $q_{2}$，其间距离为 $r$ 时，它们之间的相互作用力的大小为 $F=k\frac{q_{1}q_{2}}{r^{2}}$，式中 $k$ 为静电力常量。若用国际单位制的基本单位表示，$k$ 的单位应为\xzanswer{B}

\fourchoices[answer=B]
{kg$\cdot$A$^{2}\cdot$m$^{3}$}
{kg$\cdot$A$^{-2}\cdot$m$^{3}\cdot$s$^{-1}$}
{kg$\cdot$m$^{2}\cdot$C$^{-2}$}
{N$\cdot$m$^{2}\cdot$A$^{-2}$}

%% answer: B
%% memo:
\memoanswer{
	静电力常量 $k=9.0\times10^{9}\UNmqCq$，由牛顿第二定律知 $1\UN=1\Ukgmsq$，由电流的定义知 $1\UC=1\UAcdots$，故 $1\UNmqCq=1\UkgmcmsqAq$。选项 B 正确。
}

\item
%% number: 16
%% paperName: 2015年安徽理综第 16 题 6 分
%% typeId: 1
%% type: 单选题
%% chapter: 交流电
%% point: 变压器
%% method: 无
%% score: 6
%% degree: 450
%% duplicateId: 0
%% body:
图示电路中，变压器为理想变压器，a、b 接在电压有效值不变的交流电源两端，$R_{0}$ 为定值电阻，$R$ 为滑动变阻器。现将变阻器的滑片从一个位置滑动到另一位置，观察到电流表 A$_{1}$ 的示数增大了 0.2A，电流表 A$_{2}$ 的示数增大了 0.8A，则下列说法正确的是\xzanswer{D}

\onepicture[h=2.6cm]{figs/16.png}

\fourchoices[answer=D]
{电压表 V$_{1}$ 示数增大}
{电压表 V$_{2}$、V$_{3}$ 示数均增大}
{该变压器起升压作用}
{变阻器滑片是沿 c$\to$d 的方向滑动}

%% answer: D
%% memo:
\memoanswer{
	由于理想变压器，输入与输出端电压对确定的输入电源是一定的，A、B 错误；而 $\Delta P_{1}=\Delta P_{2}$，即 $U_{1}\Delta I_{1}=U_{2}\Delta I_{2}$，故 $\frac{\Delta I_{1}}{\Delta I_{2}}=\frac{U_{2}}{U_{1}}=\frac{1}{4}$，由变压器原理 $\frac{n_{1}}{n_{2}}=\frac{U_{1}}{U_{2}}=\frac{4}{1}$ 为降压变压器，C 错误；$V_{2}$ 不变，电流表 $A_{2}$ 的示数增大，负载电阻减小，变阻器滑片是沿 $c\to d$ 的方向滑动，D 正确。选项 D 正确。
}

\item
%% number: 17
%% paperName: 2015年安徽理综第 17 题 6 分
%% typeId: 1
%% type: 单选题
%% chapter: 电路
%% point: 电流强度
%% method: 无
%% score: 6
%% degree: 450
%% duplicateId: 0
%% body:
一根长为 $L$，横截面积为 $S$ 的金属棒，其材料的电阻率为 $\rho$，棒内单位体积自由电子数为 $n$，电子的质量为 $m$，电荷量为 $e$。在棒两端加上恒定的电压时，棒内产生电流，自由电子定向运动的平均速率为 $v$，则金属棒内的电场强度大小为\xzanswer{C}

\onepicture[h=1.6cm]{figs/17.png}

\fourchoices[answer=C]
{$\frac{mv^{2}}{2eL}$}
{$\frac{mv^{2}Sn}{e}$}
{$\rho nev$}
{$\frac{\rho ev}{SL}$}

%% answer: C
%% memo:
\memoanswer{
	由题意金属棒中的电荷量为 $q=LSne$，由一端运动到另一端所需时间 $t=L/v$，故电流 $I=q/t=Snev$，棒两端加上恒定的电压 $U=IR$，而 $R=\rho L/S$，所以 $U=\rho Lnev$；金属棒内的电场强度大小 $E=U/L=\rho nev$，选项 C 正确。
}

\item
%% number: 18
%% paperName: 2015年安徽理综第 18 题 6 分
%% typeId: 1
%% type: 单选题
%% chapter: 光学
%% point: 光的折射
%% method: 无
%% score: 6
%% degree: 450
%% duplicateId: 0
%% body:
如图所示，一束单光从空气入射到棱镜的 AB 面上，经 AB 和 AC 两个面折射后从 AC 面进入空气。当出射角 $i'$ 和入射角 $i$ 相等时，出射光线相对于入射光线偏转的角度为 $\theta$。已知棱镜顶角为 $\alpha$，则计算棱镜对该色光的折射率表达式为\xzanswer{A}

\onepicture[h=3.4cm]{figs/18.png}

\fourchoices[answer=A]
{$\frac{\sin\frac{\alpha+\theta}{2}}{\sin\frac{\alpha}{2}}$}
{$\frac{\sin\frac{\alpha+\theta}{2}}{\sin\frac{\theta}{2}}$}
{$\frac{\sin\theta}{\sin(\theta-\frac{\alpha}{2})}$}
{$\frac{\sin\theta}{\sin(\alpha-\frac{\theta}{2})}$}

%% answer: A
%% memo:
\memoanswer{
	由折射定律公式 $\frac{\sin\theta_{1}}{\sin\theta_{2}}=n$，由几何关系可得，棱镜中的折射角 $\theta_{2}=\frac{180^{\circ}-(180^{\circ}-\alpha)}{2}=\frac{\alpha}{2}$，空气中的入射角 $\theta_{1}=\frac{\alpha}{2}+\frac{\theta}{2}$，代入得 $n=\frac{\sin\frac{\alpha+\theta}{2}}{\sin\frac{\alpha}{2}}$，选项 A 正确。
}

\item
%% number: 19
%% paperName: 2015年安徽理综第 19 题 6 分
%% typeId: 1
%% type: 单选题
%% chapter: 电磁感应
%% point: 电磁感应电路分析
%% method: 等效替代
%% score: 6
%% degree: 450
%% duplicateId: 0
%% body:
如图所示，abcd 为水平放置的平行线“匚”形光滑金属导轨，间距为 $l$，导轨间有垂直于导轨平面的匀强磁场，磁感应强度大小为 $B$，导轨电阻不计。已知金属杆 MN 倾斜放置，与导轨成 $\theta$ 角，单位长度的电阻为 $r$，保持金属杆以速度 $v$ 沿平行于 cd 的方向滑动（金属杆滑动过程中与导轨接触良好）。则\xzanswer{B}

\onepicture[h=3.4cm]{figs/19.png}

\fourchoices[answer=B]
{电路中感应电动势为 $\frac{Blv}{\sin\theta}$}
{电路中感应电流的大小为 $\frac{Bv\sin\theta}{r}$}
{金属杆所受安培力的大小为 $\frac{B^{2}lv\sin\theta}{r}$}
{金属杆的热功率为 $\frac{B^{2}lv^{2}}{r\sin\theta}$}

%% answer: B
%% memo:
\memoanswer{
	由法拉第电磁感应定律可得 $E=Blv$，A 错误；又由欧姆定律 $I=\frac{E}{R}=\frac{E\sin\theta}{lr}$，故电路中感应电流的大小为 $I=\frac{Bv\sin\theta}{r}$，选项 B 正确；金属杆所受安培力的大小为 $F=BI\frac{l}{\sin\theta}=\frac{B^{2}lv}{r}$，C 错误；金属杆的热功率为 $P_{Q}=I^{2}R=\left(\frac{Bv\sin\theta}{r}\right)^{2}\frac{lr}{\sin\theta}=\frac{B^{2}lv^{2}\sin\theta}{r}$，D 错误。
}

\item
%% number: 20
%% paperName: 2015年安徽理综第 20 题 6 分
%% typeId: 1
%% type: 单选题
%% chapter: 电路
%% point: 电容
%% method: 无
%% score: 6
%% degree: 250
%% duplicateId: 0
%% body:
已知均匀带电的无穷大平面在真空中触发电场的场强大小为 $\frac{\sigma}{2\varepsilon_{0}}$，其中 $\sigma$ 为平面上单位面积所带的电荷量，$\varepsilon_{0}$ 为常量。如图所示的平行板电容器，极板正对面积为 $S$，其间为真空，带电荷量为 $Q$。不计边缘效应时，极板可看做无穷大导体板，则极板间的电场强度大小和两极板间相互的静电引力大小分别是\xzanswer{D}

\onepicture[h=3.6cm]{figs/20.png}

\fourchoices[answer=D]
{$\frac{Q}{\varepsilon_{0}S}$ 和 $\frac{Q^{2}}{\varepsilon_{0}S}$}
{$\frac{Q}{2\varepsilon_{0}S}$ 和 $\frac{Q^{2}}{\varepsilon_{0}S}$}
{$\frac{Q}{2\varepsilon_{0}S}$ 和 $\frac{Q^{2}}{2\varepsilon_{0}S}$}
{$\frac{Q}{\varepsilon_{0}S}$ 和 $\frac{Q^{2}}{2\varepsilon_{0}S}$}

%% answer: D
%% memo:
\memoanswer{
	均匀带电的无穷大平面在真空中触发电场的场强大小为 $\frac{\sigma}{2\varepsilon_{0}}$，正极板的电荷在极板间触发电场的场强大小为 $\frac{Q}{2\varepsilon_{0}S}$，负极板的电荷在极板间触发电场的场强大小为 $\frac{Q}{2\varepsilon_{0}S}$，方向相同，故极板间的电场强度大小为 $\frac{Q}{\varepsilon_{0}S}$；以一极板电荷为试探电荷得两极板间相互的静电引力为 $F=nqE=Q\cdot\frac{Q}{2\varepsilon_{0}S}$。选项 D 正确。
}

\item
%% number: 21
%% paperName: 2015年安徽理综第 21 题 12 分
%% typeId: 4
%% type: 实验题
%% chapter: 电路
%% point: 测电动势与内阻
%% method: 无
%% score: 12
%% degree: 450
%% duplicateId: 0
%% body:
（18 分）Ⅰ．在“验证力的平行四边形定则”实验中，某同学用图钉把白纸固定在水平放置的木板上，将橡皮条的一端固定在板上一点，两个细绳套系在橡皮条的另一端。用两个弹簧测力计分别拉住两个细绳套，互成角度地施加拉力，使橡皮条伸长，结点到达纸面上某一位置，如图所示。请将以下的实验操作和处理补充完整：

①用铅笔描下结点位置，记为 O；

②记录两个弹簧测力计的示数 $F_{1}$ 和 $F_{2}$，沿每条细绳（套）的方向用铅笔分别描出几个点，用刻度尺把相应的点连成线；

③只用一个弹簧测力计，通过细绳套把橡皮条的结点拉到位置 O，记录测力计的示数 $F_{3}$，\tkanswer[3.5cm]{沿此时细绳（套）的方向用铅笔描出几个点，用刻度尺把这些点连成直线}；

④按照力的图示要求，作出拉力 $F_{1}$、$F_{2}$、$F_{3}$；

⑤根据力的平行四边形定则作出 $F_{1}$ 和 $F_{2}$ 的合力 $F$；

⑥比较\tkanswer[1.2cm]{$F$ 和 $F_{3}$}的一致程度，若有较大差异，对其原因进行分析，并作出相应的改进后再次进行实验。

\onepicture[h=4.2cm, align=right]{figs/21.png}

Ⅱ．某同学为了测量一节电池的电动势和内阻，从实验室找到以下器材：一个满偏电流为 $100\UuA$、内阻为 $2500\UO$ 的表头，一个开关，两个电阻箱（$0\sim999.9\UO$）和若干导线。

（1）由于表头量程偏小，该同学首先需将表头改装成量程为 $50\UmA$ 的电流表，则应该将表头与电阻箱\tkanswer[1.2cm]{并联}（填“并联”或“串联”），并将该电阻箱阻值调为\tkanswer[1.4cm]{$5.0\UO$}。

（2）接着该同学用改装的电流表对电池的电动势及内阻进行测量，实验电路如图 1 所示，通过改变电阻 $R$ 测相应的电流 $I$，且作相关计算后一并记录如下表。

\begin{center}
\begin{tabular}{|c|c|c|c|c|c|c|}
\hline
 & 1 & 2 & 3 & 4 & 5 & 6 \\
\hline
$R$（$\UO$） & 95.0 & 75.0 & 55.0 & 45.0 & 35.0 & 25.0 \\
\hline
$I$（$\UmA$） & 15.0 & 18.7 & 24.8 & 29.5 & 36.0 & 48.0 \\
\hline
$IR$（$\UV$） & 1.42 & 1.40 & 1.36 & 1.33 & 1.26 & 1.20 \\
\hline
\end{tabular}
\end{center}

\onepicture[h=3.0cm]{figs/21a.png}

根据表中数据，图 2 中已描绘出四个点，请将第 5、6 两组数据也描绘在图 2 中，并画出 $IR-I$ 图线；

\onepicture[h=3.6cm]{figs/21b.png}

②根据图线可得电池的电动势 $E$ 是\tkanswer[1.4cm]{$1.53\UV$}，内阻 $r$ 是\tkanswer[1.4cm]{$2.0\UO$}。

%% answer:
\jdanswer{
\begin{enumerate}
	\item
	（Ⅰ）③沿此时细绳（套）的方向用铅笔描出几个点，用刻度尺把这些点连成直线；⑥$F$ 和 $F_{3}$。

	\item
	（Ⅱ）（1）并联，$5.0\UO$；（2）①描点、连线如图；②$1.53\UV$，$2.0\UO$。
\end{enumerate}
}
%% memo:
\memoanswer{
	（Ⅰ）本题原卷及材料中未提供详解，待补充。

	（Ⅱ）电流表扩大量程需并联一个电阻，有 $I_{\mathrm{mA}}R_{\mathrm{mA}}=(I-I_{\mathrm{mA}})R_{\text{并}}$，得 $R_{\text{并}}=5.0\UO$。描点作图如图，由图可得电池的电动势 $E$ 是 $1.53\UV$，曲线斜率的绝对值为电池内阻 $r$ 是 $2.0\UO$。

	\onepicture[h=3.6cm]{figs/21c.png}
}

\item
%% number: 22
%% paperName: 2015年安徽理综第 22 题 14 分
%% typeId: 6
%% type: 计算题
%% chapter: 功和能
%% point: 动能定理
%% method: 无
%% score: 14
%% degree: 650
%% duplicateId: 0
%% body:
一质量为 $0.5\Ukg$ 的小物块放在水平地面上的 A 点，距离 A 点 $5\Um$ 的位置 B 处是一面墙，如图所示。物块以 $v_{0}=9\Ums$ 的初速度从 A 点沿 AB 方向运动，在与墙壁碰撞前瞬间的速度为 $7\Ums$，碰后以 $6\Ums$ 的速度反向运动直至静止。$g$ 取 $10\Umsq$。

\begin{enumerate}
	\item
	求物块与地面间的动摩擦因数 $\mu$；

	\item
	若碰撞时间为 $0.05\Us$，求碰撞过程中墙面对物块平均作用力的大小 $F$；

	\item
	求物块在反向运动过程中克服摩擦力所做的功 $W$。
\end{enumerate}

\onepicture[h=1.8cm, align=right]{figs/22.png}

%% answer:
\jdanswer{
\begin{enumerate}
	\item
	$\mu=0.32$

	\item
	$F=130\UN$

	\item
	$W=9\UJ$
\end{enumerate}
}
%% memo:
\memoanswer{
	（1）由动能定理，有

	$-\mu mgs=\frac{1}{2}mv^{2}-\frac{1}{2}mv_{0}^{2}$，

	可得 $\mu=0.32$。

	（2）由动量定理，有 $F\Delta t=mv'-mv$，

	可得 $F=130\UN$。

	（3）$W=\frac{1}{2}mv'^{2}=9\UJ$。
}

\item
%% number: 23
%% paperName: 2015年安徽理综第 23 题 16 分
%% typeId: 6
%% type: 计算题
%% chapter: 电场
%% point: 带电粒子的偏转
%% method: 对称法
%% score: 16
%% degree: 450
%% duplicateId: 0
%% body:
在 $xOy$ 平面内，有沿 $y$ 轴负方向的匀强电场，场强大小为 $E$（图中未画出）。由 A 点斜射出一质量为 $m$，带电量为 $+q$ 的粒子，B 和 C 是粒子运动轨迹上的两点，如图所示。其中 $l_{0}$ 为常数。粒子所受重力忽略不计。求：

\begin{enumerate}
	\item
	粒子从 A 到 C 过程中电场力对它做的功；

	\item
	粒子从 A 到 C 过程所经历的时间；

	\item
	粒子经过 C 点时的速率。
\end{enumerate}

\onepicture[h=4.0cm, align=right]{figs/23.png}

%% answer:
\jdanswer{
\begin{enumerate}
	\item
	$W_{AC}=3qEl_{0}$

	\item
	$t=3\sqrt{\frac{2ml_{0}}{qE}}$

	\item
	$v_{C}=\sqrt{\frac{17qEl_{0}}{2m}}$
\end{enumerate}
}
%% memo:
\memoanswer{
	（1）$W_{AC}=qE(y_{A}-y_{C})=3qEl_{0}$。

	（2）根据抛体运动的特点，粒子在 $x$ 方向上做匀速直线运动，由对称性可知轨道最高点 D 在 $y$ 轴上，可令 $t_{AD}=t_{DB}=T$，则 $t_{BC}=T$。

	由 $qE=ma$ 得 $a=\frac{qE}{m}$，

	又 $y_{D}=\frac{1}{2}aT^{2}$，$y_{D}+3l_{0}=\frac{1}{2}a(2T)^{2}$，

	解得 $T=\sqrt{\frac{2ml_{0}}{qE}}$，

	则 A→C 过程所经历的时间 $t=3\sqrt{\frac{2ml_{0}}{qE}}$。

	（3）粒子在 DC 段做类平抛运动，于是有

	$2l_{0}=v_{Cx}(2T)$，$v_{Cy}=a(2T)$，

	$v_{C}=\sqrt{v_{Cx}^{2}+v_{Cy}^{2}}=\sqrt{\frac{17qEl_{0}}{2m}}$。
}

\item
%% number: 24
%% paperName: 2015年安徽理综第 24 题 20 分
%% typeId: 6
%% type: 计算题
%% chapter: 曲线运动
%% point: 双星问题
%% method: 无
%% score: 20
%% degree: 250
%% duplicateId: 0
%% body:
由三颗星体构成的系统，忽略其他星体对它们的作用，存在着一种运动形式：三颗星体在相互之间的万有引力作用下，分别位于等边三角形的三个顶点上，绕某一共同的圆心 O 在三角形所在平面内做相同角速度的圆周运动（图示为 A、B、C 三颗星体质量不相同时的一般情况），若 A 星体质量为 $2m$，B、C 两星体的质量均为 $m$，三角形的边长为 $a$，求：

\begin{enumerate}
	\item
	A 星体所受合力大小 $F_{A}$；

	\item
	B 星体所受合力大小 $F_{B}$；

	\item
	C 星体的轨道半径 $R_{C}$；

	\item
	三星体做圆周运动的周期 $T$。
\end{enumerate}

\onepicture[h=4.2cm, align=right]{figs/24.png}

%% answer:
\jdanswer{
\begin{enumerate}
	\item
	$F_{A}=2\sqrt{3}G\frac{m^{2}}{a^{2}}$

	\item
	$F_{B}=\sqrt{7}G\frac{m^{2}}{a^{2}}$

	\item
	$R_{C}=\frac{\sqrt{7}}{4}a$

	\item
	$T=\pi\sqrt{\frac{a^{3}}{Gm}}$
\end{enumerate}
}
%% memo:
\memoanswer{
	（1）由万有引力定律，A 星体所受 B、C 星体引力大小为

	$F_{BA}=G\frac{m_{A}m_{B}}{r^{2}}=G\frac{2m^{2}}{a^{2}}=F_{CA}$，方向如图，

	则合力大小为 $F_{A}=2\sqrt{3}G\frac{m^{2}}{a^{2}}$。

	（2）同上，B 星体所受 A、C 星体引力大小分别为

	$F_{AB}=G\frac{m_{A}m_{B}}{r^{2}}=G\frac{2m^{2}}{a^{2}}$，$F_{CB}=G\frac{m_{c}m_{B}}{r^{2}}=G\frac{m^{2}}{a^{2}}$，方向如图，

	由 $F_{Bx}=F_{AB}\cos60^{\circ}+F_{CB}=2G\frac{m^{2}}{a^{2}}$，$F_{By}=F_{AB}\sin60^{\circ}=\sqrt{3}G\frac{m^{2}}{a^{2}}$，

	可得 $F_{B}=\sqrt{F_{Bx}^{2}+F_{By}^{2}}=\sqrt{7}G\frac{m^{2}}{a^{2}}$。

	（3）通过分析可知，圆心 O 在中垂线 AD 的中点，$R_{C}=\sqrt{\left(\frac{\sqrt{3}}{4}a\right)^{2}+\left(\frac{1}{2}a\right)^{2}}$，

	（或：由对称性可知 $OB=OC=R_{C}$，$\cos\angle OBD=\frac{F_{Bx}}{F_{B}}=\frac{DB}{OB}=\frac{\frac{1}{2}a}{R_{C}}$，

	可得 $R_{C}=\frac{\sqrt{7}}{4}a$。

	（4）三星体运动周期相同，对 C 星体，由 $F_{C}=F_{B}=\sqrt{7}G\frac{m^{2}}{a^{2}}=m\left(\frac{2\pi}{T}\right)^{2}R_{C}$，

	可得 $T=\pi\sqrt{\frac{a^{3}}{Gm}}$。

	\onepicture[h=4.2cm]{figs/24b.png}
}

\end{enumerate}

\end{document}
